\label{sec:introduction}

The decennial redistricting cycle is not an event---it is a process that begins
the moment the prior census data arrives and extends through years of litigation
over the resulting maps. The planning literature on redistricting has focused
almost entirely on the process \textit{after} census data release: how maps are
drawn, what criteria apply, who decides, and how courts review the outcome.
What the literature has largely neglected is the preparation process \textit{before}
census data release: the infrastructure, data, and algorithmic foundations that
must be in place before the redistricting window opens.

For human mapmakers, this preparation neglect is partially defensible. Human mappers
draw maps using the census data itself; the main prerequisite is software training
and political strategy, neither of which requires years of advance preparation.
For algorithmic redistricting systems like \textsc{bisect}, the situation is categorically
different. Algorithmic redistricting requires:

\begin{enumerate}
\item \textbf{Geographic adjacency graphs} built from the census geography of the
  target year. These graphs encode which census units share borders, and they take
  weeks to months to construct at the block-group level.
\item \textbf{Validated demographic data} for every census unit in every state,
  including total population, minority VAP, and commuting flows. The data validation
  process catches geographic anomalies (water tracts, zero-population tracts,
  prison population allocation) that cause algorithmic errors.
\item \textbf{Tested algorithm configuration} specifying the structure, weights,
  and search parameters for each state's redistricting. Changes to this configuration
  after census data release trigger legal transparency challenges.
\item \textbf{Computational infrastructure} capable of running 50-state sweeps within
  the redistricting window. At \textsc{bisect}'s current performance of $\sim$90 minutes
  for a 50-state tract-resolution sweep, the block-group resolution sweep estimated
  for 2030 will require $\sim$90 minutes as well, but with substantially more
  memory requirements.
\end{enumerate}

None of these prerequisites can be assembled in weeks. The adjacency graph construction
alone requires months of computation and manual validation. The implication is direct:
if redistricting infrastructure is not substantially complete before the 2031 census
data release, the 6--12 month redistricting window will be consumed by infrastructure
construction rather than redistricting itself.

\subsection{The Current State of \textsc{bisect} Infrastructure}

As of 2026, the \textsc{bisect} platform has:
\begin{itemize}
\item Tract-level adjacency graphs for all 50 states, validated for 2000, 2010, and 2020.
\item 50-state congressional redistricting capability at tract resolution.
\item Block-group adjacency graphs for zero states (not yet constructed).
\item 2020 census data fully loaded and validated.
\item LODES 2021 commuting data and ACS 2019--2021 5-year data loaded.
\end{itemize}

The infrastructure gap for 2030 is therefore primarily the block-group adjacency graphs
and the updated demographic data. This paper provides the timeline and resource estimates
for closing this gap before the 2031 redistricting window opens.

\subsection{Organization}

Section~\ref{sec:timeline} describes the 2030 redistricting timeline, including the
key statutory deadlines and the uncertainty ranges around census data release.
Section~\ref{sec:preparation} describes the four preparation tasks in detail.
Section~\ref{sec:dia-seed} explains the DIA statutory seed formula and why it
prohibits last-minute algorithm changes. Section~\ref{sec:conclusion} provides
the roadmap for Q.1--Q.4 and the recommended preparation schedule.
